%%%%%%%%%%%%%%%%%%%%%%%%%%%%%%%%%%%%%%%%%%%%%%%%%%%%%%%
% TikZ Modular Mechanical Components Library
% Description: Scalable and reusable blocks for mechanical systems
%%%%%%%%%%%%%%%%%%%%%%%%%%%%%%%%%%%%%%%%%%%%%%%%%%%%%%%

% === SPRING-DAMPER BLOCK ===
\newcommand{\springDamperBlock}[3]{%
  % #1 = anchor point (e.g., (0,0))
  % #2 = index for labels (e.g., 1 → k_1, c_1)
  % #3 = TikZ options (e.g., rotate=90, scale=0.8)
  \begin{scope}[shift={#1}, #3]0 +6 
    \coordinate (A) at (0,0.5);
    \coordinate (B) at (1,0.5);
    \coordinate (C) at (0,0);
    \coordinate (D) at (1,0);

    % Vertical connectors
    \draw[thick] (A) -- (C);
    \draw[thick] (B) -- (D);

    % Spring
    \draw[decorate,decoration={coil,aspect=0.5,segment length=2mm,amplitude=2mm}] (A) -- (B);
    \node at (0.5,1) {\( k_{#2} \)};

    % Damper
    \draw[thick] (0,0) -- (0.45,0);
    \draw[thick] (0.45,0.2) -- (0.45,-0.2);
    \draw[thick] (0.45,0.2) -- (0.6,0.2);
    \draw[thick] (0.45,-0.2) -- (0.6,-0.2);
    \draw[thick] (0.55,0) -- (1,0);
    \draw[thick] (0.55,0.15) -- (0.55,-0.15);
    \node at (0.5,-0.5) {\( c_{#2} \)};
  \end{scope}
}


% === WALL BLOCK ===
\newcommand{\wallBlock}[4]{%
  % #1 = anchor point
  % #2 = length of the wall
  % #3 = TikZ options
  % #4 = scale factor for step size

  \begin{scope}[shift={#1}, #3]
    \draw[thick] (0,0) -- (#2,0);

    \pgfmathsetmacro{\s}{#4}
    \pgfmathsetmacro{\step}{0.2*\s}
    \pgfmathsetmacro{\n}{#2}
    \pgfmathsetmacro{\N}{int(\n/\step)}

    \foreach \i in {0,...,\N} {
      \pgfmathsetmacro{\x}{\i*\step}
      \draw (\x,0) -- (\x-\step,-\step);
    }
  \end{scope}
}

% === MASS ON WHEELS ===
\newcommand{\massOnWheels}[4]{%
  % #1 = anchor point
  % #2 = scale factor
  % #3 = label index (e.g., 1 → m_1)
  % #4 = TikZ options
  \begin{scope}[shift={#1}, #4]
    \pgfmathsetmacro{\s}{#2}
    \pgfmathsetmacro{\h}{1.0*\s}
    \pgfmathsetmacro{\w}{2.5*\s}
    \pgfmathsetmacro{\r}{0.2*\s}

    \draw[thick, fill=gray!10] (0,0) rectangle (\w,\h);
    \node at (\w/2,\h/2) {\( m_{#3}, J_{#3} \)};

    \draw[thick] ({\r + 0.15*\s}, -\r) circle(\r);
    \draw[thick] ({\w - \r - 0.15*\s}, -\r) circle(\r);

    % Draw floor
    \wallBlock{({0-2*\r},{0-2*\r})}{\w+4*\r}{#4}{#2}
  \end{scope}
}

% === MASS BLOCK ===
\newcommand{\massBlock}[4]{%
  % #1 = anchor point
  % #2 = scale factor
  % #3 = label index (e.g., 2 → m_2)
  % #4 = TikZ options
  \begin{scope}[shift={#1}, #4]
    \pgfmathsetmacro{\s}{#2}
    \pgfmathsetmacro{\h}{1.0*\s}
    \pgfmathsetmacro{\w}{2.5*\s}

    \draw[thick, fill=gray!10] (0,0) rectangle (\w,\h);
    \node at (\w/2,\h/2) {\( m_{#3} \)};
  \end{scope}
}

% === DISK ===
\newcommand{\disk}[4]{%
  % #1 = center point
  % #2 = radius
  % #3 = label index
  % #4 = TikZ options
  \begin{scope}[shift={#1}, #4]
    \pgfmathsetmacro{\r}{#2}

    \draw[thick] (0,0) circle (\r);
    \draw[thick] (0,0) circle (0.1*\r);

    \node[anchor=west] at (-\r*1.625,\r) {\( m_{#3}, J_{#3} \)};

    \draw[->, thick, blue] (0,0) -- ({\r*cos(30)},{\r*sin(30)});
    \node[blue] at ({\r*cos(30)-0.7*\r},{\r*sin(30)}) {\( R_{#3} \)};
  \end{scope}
}

% === FIXED DISK ===
\newcommand{\fixedDisk}[4]{%
  % #1 = center point
  % #2 = radius
  % #3 = label index
  % #4 = TikZ options
  \begin{scope}[shift={#1}, #4]
    \pgfmathsetmacro{\r}{#2}

    \draw[thick] (0,0) circle (\r);
    \draw[thick] (0,0) circle (0.1*\r);
    \draw[thick] (-0.1*\r,-0.1*\r) -- ({0.2*\r*cos(230)}, {0.2*\r*sin(230)});
    \draw[thick] (0,-0.12*\r) -- ({0.22*\r*cos(270)}, {0.2*\r*sin(270)});
    \draw[thick] (+0.1*\r,-0.1*\r) -- ({0.2*\r*cos(310)}, {0.2*\r*sin(310)});

    \node[anchor=west] at ({\r*0.625},\r) {\( m_{#3}, J_{#3} \)};

    \draw[->, thick, blue] (0,0) -- ({\r*cos(30)},{\r*sin(30)});
    \node[blue] at ({\r*cos(30)-0.7*\r},{\r*sin(30)}) {\( R_{#3} \)};
  \end{scope}
}

% === DOUBLE WELDED DISK ===
\newcommand{\doubleWeldedDisk}[4]{%
  % #1 = center point
  % #2 = inner radius (Rb)
  % #3 = label index
  % #4 = TikZ options
  \begin{scope}[shift={#1}, #4]
    \pgfmathsetmacro{\Rb}{#2}
    \pgfmathsetmacro{\Ra}{1.6*\Rb}

    \draw[line width=1.5pt] (0,0) circle (\Ra);
    \draw[thick] (0,0) circle (\Rb);

    \node[anchor=west] at (\Rb,\Ra) {\( m_{#3}, J_{#3} \)};

    \draw[->, thick, blue] (0,0) -- ({\Ra*cos(30)},{\Ra*sin(30)});
    \node[blue] at ({\Ra*cos(30)-0.35*\Ra},{\Ra*sin(30)+0.1*\Ra}) {\( R_{#3a} \)};

    \draw[->, thick, blue] (0,0) -- ({\Rb*cos(-45)},{\Rb*sin(-45)});
    \node[blue] at ({\Rb*cos(-45)-0.6*\Rb},{\Rb*sin(-45)+0.1*\Rb}) {\( R_{#3b} \)};
  \end{scope}
}

% === CONVENTIONS BLOCK ===
\newcommand{\conventionsBlock}[2]{%
  % #1 = anchor point (e.g., (0,0))
  % #2 = TikZ options (e.g., rotate=90, scale=0.8)
  \begin{scope}[shift={#1}, #2]

    % Axes (x and y)
    \draw[->, thick, gray] (0,0) -- (0,2);
    \draw[->, thick, gray] (0,0) -- (2,0);
    % Axis labels
    \node[text=gray] at (2,-0.2) {\( \hat{x} \)};
    \node[text=gray] at (-0.2,2) {\( \hat{y} \)};

    % Positive rotation (circular arrow)
    \draw[->, thick, gray] (1.2,0) arc[start angle=0, end angle=90, radius=1.2];
    %label
    \node[text=gray] at (1,1) {\( \hat{\theta} \)};

    % Sign convention
    \node[text=gray] at (0.5,0.5) {\( + \)};

    % Spring horizontal (symbolic)
    \draw[->, thick, gray] (0.5,-0.7) -- (0,-0.7);
    \draw[decorate, decoration={coil,aspect=0.5,segment length=2mm,amplitude=2mm}, gray]
      (0.5,-0.7) -- (1.5,-0.7);
    \draw[->, thick, gray] (1.5,-0.7) -- (2,-0.7);
    \node[text=gray] at (1,-0.3) {\( + \)};

  \end{scope}
}

% === ROLLER ===
\newcommand{\roller}[3]{
    % #1 = anchor point (e.g., (0,0))
    % #2 = scale factor
    % #3 = TikZ options (e.g., rotate)

    \begin{scope}[shift={#1}, #3]
        \pgfmathsetmacro{\s}{#2}
        \pgfmathsetmacro{\h}{1.0*\s}
        \pgfmathsetmacro{\b}{1.4*\s}
        \pgfmathsetmacro{\r}{0.15*\s}
        \pgfmathsetmacro{\f}{1.7*\s}

        % Coordinates relative to (0,0)
        \coordinate (A) at ({-\b/2}, {-\h});
        \coordinate (B) at ({\b/2}, {-\h});
        \coordinate (C) at (0,0);
        \coordinate (D) at ({-\b/2+2*\r}, {-\h - \r});
        \coordinate (E) at ({\b/2-2*\r}, {-\h - \r});
        \coordinate (F) at ({-\b/2 - \r}, {-\h - 2*\r});

        % Draw triangle (roller body)
        \draw[thick] (A) -- (B) -- (C) -- cycle;

        % Draw wheels
        \draw[thick] (D) circle(\r);
        \draw[thick] (E) circle(\r);

        % Draw floor
        \wallBlock{(F)}{\f}{#3}{#2}
    \end{scope}
}

% === DISK ===
\newcommand{\diskRoller}[4]{%
  % #1 = center point
  % #2 = radius
  % #3 = label index
  % #4 = TikZ options
  \begin{scope}[shift={#1}, #4]
    \pgfmathsetmacro{\r}{#2}

    \draw[thick] (0,0) circle (\r);

    \node[anchor=west] at ({\r*0.625},\r) {\( m_{#3}, J_{#3} \)};

    \draw[->, thick, blue] (0,0) -- ({\r*cos(30)},{\r*sin(30)});
    \node[blue] at ({\r*cos(30)-0.7*\r},{\r*sin(30)}) {\( R_{#3} \)};

    \roller{(0,0)}{0.3*\r}{#4}
  \end{scope}
}